%% secoes/01-introducao.tex — incluído por main.tex

%%=====================================================================%%
\section{Introdução}\label{sec:introducao}
%%=====================================================================%%

O período neonatal concentra a maior parcela da mortalidade infantil \cite{unigme2024}. 
No ano de 2023, 4,8 milhões de crianças morreram antes dos cinco anos no mundo, sendo 2,3 milhões delas recém-nascidos \cite{unigme2024}. 
Mais do que um período de risco elevado, é a janela em que se assentam as bases da sobrevivência e do desenvolvimento, de modo que as condições em que uma criança nasce repercutem ao longo de toda a sua vida \cite{lawn2014}. Nesse cenário, os desfechos neonatais adversos, eventos que sinalizam comprometimento da saúde já no nascimento, funcionam como indicadores sensíveis de saúde \cite{lawn2014}.
Este trabalho trata de quatro desfechos neonatais adversos: baixo peso ao nascer, nascimento pré-termo, malformação congênita e Apgar baixo no quinto minuto.

%Esses desfechos não se distribuem ao acaso: quanto menor a posição socioeconômica, pior o estado de saúde, padrão que se reproduz tanto entre países quanto dentro deles \cite{marmot2008}. O Brasil é um caso emblemático dessa tensão. Apesar dos avanços expressivos na saúde materno-infantil nas últimas décadas, impulsionados por crescimento econômico, redução das desigualdades de renda, expansão da atenção básica e programas de transferência de renda \cite{victora2011}, persistem profundas iniquidades entre regiões e grupos sociais, que fazem com que o risco de um desfecho adverso ao nascer dependa, ainda hoje, de onde e de quem a criança nasce.

 
%As subseções a seguir apresentam cada um deles.

%%---------------------------------------------------------------------%%
%% Subseções da introdução, um arquivo por desfecho.                   %%
%% Caminhos relativos à pasta do manuscrito, não à deste arquivo.      %%
%%---------------------------------------------------------------------%%

%% secoes/introducao/01-baixo-peso.tex — incluído por main.tex

%\subsection{Baixo peso ao nascer}\label{subsec:intro-bpn}

Define-se baixo peso ao nascer como o peso inferior a 2.500 gramas ao nascimento, independentemente da idade gestacional \cite{blencowe2019}.
O baixo peso ao nascer constitui um dos mais importantes problemas de saúde pública perinatal, por refletir as condições em que a gestação transcorreu e por anteceder desfechos adversos que se estendem muito além do nascimento \cite{falcao2020,kramer1987}. Suas causas são multifatoriais e fortemente ligadas à vulnerabilidade social, o que o torna, em grande medida, evitável por meio de cuidados adequados ao longo da gestação \cite{kramer1987,fonseca2014}. Ainda assim, apesar dos avanços em saúde materno-infantil, sua ocorrência persiste elevada, no mundo e no Brasil, sem redução expressiva nas últimas décadas, configurando uma iniquidade que recai de forma desproporcional sobre as populações socioeconomicamente mais vulneráveis \cite{falcao2020,blencowe2019}.

%Define-se baixo peso ao nascer como o peso inferior a 2.500 gramas ao nascimento, independentemente da idade gestacional \cite{blencowe2019}. Suas repercussões não se restringem ao período neonatal: além do maior risco de morbidade e mortalidade nos primeiros anos de vida, indivíduos que nascem com baixo peso apresentam, na vida adulta, maior predisposição a doenças crônicas não transmissíveis, como diabetes, hipertensão e doenças cardiovasculares, associação sustentada pela hipótese das origens desenvolvimentistas da saúde e da doença \cite{barker2007}. Trata-se, portanto, de um agravo cujas consequências acompanham o indivíduo ao longo de todo o curso de vida.

%Do ponto de vista biológico, o baixo peso ao nascer decorre de prematuridade, de restrição de crescimento intrauterino ou da combinação de ambos \cite{kramer1987}. Esses mecanismos, por sua vez, são condicionados por um amplo conjunto de determinantes: características maternas, como idade e história obstétrica; a assistência pré-natal, mensurada pelo número de consultas; e fatores socioeconômicos, como escolaridade materna, renda, condições domiciliares e região de residência \cite{falcao2020}. Essa cadeia de determinação evidencia por que o desfecho se concentra justamente nos contextos de maior privação social \cite{falcao2020}.

%Por não existir medida curativa após o nascimento, o enfrentamento do baixo peso ao nascer é essencialmente preventivo. A assistência pré-natal adequada, com número suficiente de consultas e início precoce, figura entre as estratégias mais efetivas, ao permitir o acompanhamento nutricional da gestante, a identificação oportuna de complicações e a orientação quanto a comportamentos de risco \cite{falcao2020,fonseca2014}. A esse cuidado somam-se políticas de proteção social, como os programas de transferência de renda, que atuam sobre as condições socioeconômicas subjacentes ao problema \cite{falcao2020,rasella2013}.

%Quanto à magnitude, estimou-se que, em 2015, 14,6\% dos nascimentos no mundo, cerca de 20,5 milhões, ocorreram com baixo peso, dos quais 91\% em países de baixa e média renda \cite{blencowe2019}. No Brasil, a prevalência foi estimada em aproximadamente 8,5\% em 2017, patamar próximo ao observado na América Latina e no Caribe e sem tendência de queda relevante nas últimas décadas \cite{blencowe2019,ms2017}. 
%Chama atenção, no cenário nacional, o paradoxo epidemiológico do baixo peso ao nascer, em que as maiores proporções recaem sobre as regiões Sul e Sudeste, economicamente mais desenvolvidas, evidência de que a determinação do desfecho não se reduz à dimensão econômica \cite{silva2010}.

%% secoes/introducao/02-pre-termo.tex — incluído por main.tex

%\subsection{Nascimento pré-termo}\label{subsec:intro-pt}

Já o nascimento pré-termo é definido quando ocorre antes de 37 semanas completas de gestação. A prematuridade é uma das principais causas de morte de crianças menores de cinco anos e está associada a efeitos de longo prazo sobre a saúde física, o neurodesenvolvimento e as condições socioeconômicas dos que sobrevivem \cite{ohuma2023}. Estima-se que 13,4 milhões de crianças nasceram pré-termo no mundo em 2020, 9,9\% de todos os nascimentos, sem mudança mensurável em relação a 2010 \cite{ohuma2023}.

%No Brasil, a prevalência de nascimentos pré-termo foi de 11,1\% entre 2011 e 2021 \cite{alberton2023} e passou de 11,3\% em 2014 para 11,9\% em 2023 \cite{victor2025}, acima da média mundial. As maiores prevalências recaem sobre mulheres indígenas e pretas, com menor escolaridade, nas idades extremas e com poucas consultas de pré-natal, e sobre as regiões Norte e Nordeste \cite{alberton2023,victor2025}.

%Na população de baixa renda, o número reduzido de consultas de pré-natal está entre os fatores de maior risco, tanto para o primeiro nascimento pré-termo quanto para a sua recorrência \cite{rocha2022}. As políticas de proteção social também se relacionam ao desfecho: entre mulheres de baixa renda, o recebimento do Bolsa Família durante a gestação esteve associado à redução dos nascimentos pré-termo extremos \cite{ortelan2024}.

%% secoes/introducao/03-malformacao.tex — incluído por main.tex

%\subsection{Malformação congênita}\label{subsec:intro-malf}

As malformações congênitas, ou anomalias congênitas, são alterações estruturais ou funcionais que se originam durante a vida intrauterina e podem ser identificadas antes do nascimento, ao nascer ou mais tarde na infância \cite{who2023congenital}. Estima-se que cerca de 240 mil recém-nascidos morram por ano no mundo, nos primeiros 28 dias de vida, em decorrência delas, além de outras 170 mil mortes entre um mês e cinco anos de idade \cite{who2023congenital}. A maioria não tem causa conhecida, e as anomalias graves concentram-se nos países de baixa e média renda \cite{who2023congenital}.
No Brasil, a presença de anomalia congênita ao nascer foi de 8,0 por mil nascidos vivos, valor que os autores consideram afetado pela detecção pré-natal limitada e pela subnotificação \cite{melkan2025}.

%No Brasil, a presença de anomalia congênita é registrada na declaração de nascido vivo. Entre 2011 e 2020, a prevalência ao nascer foi de 8,0 por mil nascidos vivos, valor que os autores consideram afetado pela detecção pré-natal limitada e pela subnotificação \cite{melkan2025}. Ainda que fatores genéticos e biológicos tenham papel central, há associação com condições sociais: em âmbito nacional, a chance de anomalia congênita foi maior entre filhos de mães pretas, sem consultas de pré-natal, com mais de 40 anos e em gestações múltiplas \cite{fernandes2025}. Na epidemia de Zika, os casos da síndrome congênita concentraram-se entre as mulheres socialmente mais vulneráveis \cite{paixao2022}.

%% secoes/introducao/04-apgar.tex — incluído por main.tex

%\subsection{Apgar baixo no quinto minuto}\label{subsec:intro-apgar}

Por fim, o índice de Apgar, proposto por Virginia Apgar em 1953, avalia a condição do recém-nascido logo após o parto a partir de cinco sinais: frequência cardíaca, esforço respiratório, tônus muscular, irritabilidade reflexa e cor da pele, cada um pontuado de 0 a 2 \cite{apgar1953}. A avaliação é feita no primeiro e no quinto minuto de vida, e o índice inferior a 7 no quinto minuto é usado como indicador de comprometimento das condições de nascimento e da necessidade de cuidados especializados \cite{magalhaes2023}.
No Brasil, a proporção de recém-nascidos com Apgar menor que 7 no quinto minuto caiu de 2,1\% em 1999 para 0,9\% em 2018--2019 \cite{magalhaes2023}. 
Entre os quatro desfechos deste trabalho, o Apgar baixo é o menos investigado sob a perspectiva social (Seção~\ref{sec:relacionados}).

%Prematuridade, baixo peso ao nascer e anomalia congênita são os fatores de maior impacto, mas também se associam ao desfecho características maternas e da assistência, como idade a partir de 35 anos, escolaridade inferior a oito anos, ausência de companheiro e menos de seis consultas de pré-natal \cite{magalhaes2023}. Apesar disso, entre os quatro desfechos deste trabalho, o Apgar baixo é o menos investigado sob a perspectiva social (Seção~\ref{sec:relacionados}).


%%---------------------------------------------------------------------%%
%\subsection{Objetivo e organização do trabalho}\label{subsec:intro-objetivo}
%%---------------------------------------------------------------------%%

O objetivo deste trabalho é identificar, por meio de descoberta de subgrupos, uma técnica de mineração de dados, as combinações de características associadas a maior ocorrência dos quatro desfechos em populações socioeconomicamente vulneráveis do Brasil. Para isso, foi utilizada a Coorte de 100 Milhões de Brasileiros vinculada ao SINASC, base que reúne, para famílias de baixa renda inscritas no Cadastro Único, os dados do nascimento, da mãe, do domicílio e do recebimento do Bolsa Família \cite{pita2018}. Após o tratamento dos dados, a base analisada ficou com 5.242.215 recém-nascidos entre 2011 e 2015, descritos por 42 atributos. Sobre ela, aplicou-se o SSDP+, um algoritmo evolutivo para essa tarefa \cite{lucas2018}, que retorna regras legíveis, como ``1 a 3 consultas de pré-natal e primeira gestação'', que descrevem os grupos em que o desfecho é mais frequente. A busca foi feita para cada um dos quatro desfechos em oito populações: o recorte geral, usado como referência, e sete segmentos vulneráveis (municípios de até 10 mil habitantes, população rural, indígena e quilombola, beneficiárias do Programa Bolsa Família, mães com baixa escolaridade e mães com menos de 15 anos), totalizando 32 análises. Os subgrupos encontrados foram ainda disponibilizados em um painel interativo de visualização.

%Para isso, aplica-se o algoritmo SSDP+ à Coorte de 100 Milhões de Brasileiros vinculada ao SINASC. A análise abrange o recorte geral, como referência, e sete segmentos vulneráveis: municípios de até 10 mil habitantes, população rural, indígena e quilombola, beneficiárias do Programa Bolsa Família, mães com baixa escolaridade e mães com menos de 15 anos. Busca-se, assim, evidenciar padrões de risco próprios de cada grupo, capazes de subsidiar políticas mais direcionadas. Os subgrupos encontrados foram também disponibilizados em um painel interativo de visualização.
